Una empresa cuenta con 4 plantas de producción (Orígenes: $O_1, O_2, O_3, O_4$) y debe abastecer a 4 centros de distribución (Destinos: $D_1, D_2, D_3, D_4$). En la Tabla \ref{tab:enunciado} se presentan los costes unitarios de transporte, las capacidades de oferta de cada origen y los requerimientos de demanda de cada destino.

\begin{table}[htbp]
\centering
\caption{Matriz inicial de costos, oferta y demanda.}
\label{tab:enunciado}
\begin{tabular}{cccccc}
\toprule
& $D_1$ & $D_2$ & $D_3$ & $D_4$ & \textbf{Oferta} \\
\midrule
$O_1$ & 10 & 20 & 5  & 7  & \textbf{15} \\
$O_2$ & 13 & 9  & 12 & 8  & \textbf{25} \\
$O_3$ & 4  & 15 & 7  & 9  & \textbf{20} \\
$O_4$ & 14 & 7  & 1  & 0  & \textbf{40} \\
\midrule
\textbf{Demanda} & \textbf{30} & \textbf{20} & \textbf{35} & \textbf{15} &  \\
\bottomrule
\end{tabular}
\end{table}

Resuelve hasta el óptimo este problema utilizando: 1) el método de Vogel para obtener una solución inicial y, 2) el método de MODI para calcular los costes reducidos.
